% =====================================================================
% Tracklet Splitter Evaluation
% Remaining TODOs (require additional experiments):
%   TODO:HOTA:<name>  — end-to-end HOTA with XGBoost merger (finetuning pending)
%   TODO:REID:*       — spatio-temporal ReID splitter results (separate eval)
% =====================================================================

\section{Tracklet Splitter Evaluation}\label{sec:splitter_eval}

This section evaluates the attribute-based tracklet splitter and compares it
against three alternative splitter designs explored during development: the
bounding-box anomaly splitter, the trajectory splitter, and the GTA-style
DBSCAN splitter. The spatio-temporal ReID splitter is evaluated separately in
Section~\ref{subsec:reid_splitter_eval}.

% ---------------------------------------------------------------------
\subsection{Evaluation Methodology}\label{subsec:splitter_methodology}
% ---------------------------------------------------------------------

Evaluating a tracklet splitter in isolation is non-trivial because splitting
a tracklet always increases the fragment count and therefore decreases HOTA
directly, regardless of split quality. Reporting HOTA of the splitter output
alone would conflate two distinct effects: the splitter making bad decisions,
and the evaluation metric penalising fragmentation that will be corrected by
the downstream merger. To separate these effects, two complementary evaluation
strategies are used.

\paragraph{Intrinsic metrics.}
Intrinsic metrics measure split-decision quality against the ground-truth
identity labels without any downstream component. The ground-truth identity
switches are defined as all positions inside an input tracklet where the
underlying ground-truth identity changes between two consecutive detections.
The baseline tracklets contain \textbf{1\,048} such switches across the 49
test sequences.

Three metrics are reported. \textbf{IDS-Precision} (IDS-P) is the fraction
of predicted splits that fall within $\pm\tau_f$ frames of a ground-truth
switch; a low value indicates over-splitting of clean tracklets.
\textbf{IDS-Recall} (IDS-R) is the fraction of ground-truth switches that
are detected within $\pm\tau_f$ frames; a low value indicates missed
switches. \textbf{IDS-F1} is the harmonic mean of the two, providing a
single summary of split-decision quality. A tolerance of $\tau_f = 10$
frames is used throughout. The number of output tracklets is also reported
as a measure of fragmentation cost.

Tracklet purity, defined as the mean fraction of frames in each output
tracklet that belong to its dominant ground-truth identity, is not reported
in the main table. Because the majority of tracklets are not split by any
method, the aggregate purity values cluster tightly between 0.97 and 0.98
across all configurations and do not meaningfully differentiate the methods.

\paragraph{Extrinsic evaluation.}
Each configuration is also evaluated end-to-end through the full ALERT
pipeline with a fixed downstream XGBoost merger (Section~\ref{xgboost_merger}),
reporting HOTA on the test split. Because the same merger is used for every
configuration, differences in HOTA can be primarily attributed to the splitter.

% ---------------------------------------------------------------------
\subsection{Individual Signal Analysis}\label{subsec:individual_signals}
% ---------------------------------------------------------------------

Table~\ref{tab:splitter_comparison} compares the individual splitter signals
and their combinations on the test split intrinsic metrics.

\begin{table}[H]
\centering
\begin{tabular}{lccccc}
\toprule
\textbf{Method} & \textbf{\#Tracklets} & \textbf{IDS-P} ($\uparrow$) & \textbf{IDS-R} ($\uparrow$) & \textbf{IDS-F1} ($\uparrow$) & \textbf{HOTA} ($\uparrow$) \\
\midrule
Baseline (no splitting)     & 1725 & ---    & 0.000  & ---    & \texttt{TODO:HOTA:baseline} \\
\midrule
Jersey                      & 1824 & 0.7778 & 0.0735 & 0.1343 & \texttt{TODO:HOTA:jersey} \\
Team                        & 1871 & 0.6351 & 0.0897 & 0.1572 & \texttt{TODO:HOTA:team} \\
Bounding-box anomaly        & 1875 & 0.6133 & 0.0878 & 0.1536 & \texttt{TODO:HOTA:bbox} \\
Trajectory                  & 1960 & 0.3489 & 0.0782 & 0.1278 & \texttt{TODO:HOTA:traj} \\
GTA DBSCAN                  & 2367 & 0.0755 & 0.1479 & 0.1000 & \texttt{TODO:HOTA:gta} \\
\midrule
Jersey + Team               & 1877 & 0.6234 & 0.0916 & 0.1597 & \texttt{TODO:HOTA:jersey+team} \\
Jersey + Team + Bbox        & 1930 & 0.5266 & 0.1040 & 0.1737 & \texttt{TODO:HOTA:jersey+team+bbox} \\
Jersey + Team + Bbox + Traj & 2072 & 0.3324 & 0.1107 & 0.1661 & \texttt{TODO:HOTA:jersey+team+bbox+traj} \\
\bottomrule
\end{tabular}
\caption{Intrinsic and extrinsic performance of tracklet splitter configurations
on the 49 test sequences. All configurations are evaluated against the same
1\,048 ground-truth identity switches in the baseline Deep-EIoU tracklets.
Extrinsic HOTA is computed end-to-end with the XGBoost merger as a fixed
downstream component.}
\label{tab:splitter_comparison}
\end{table}

\paragraph{Jersey splitter.}
The jersey-number splitter achieves the highest precision of all individual
splitters (77.78\%), meaning that nearly four in five predicted splits
correspond to a genuine identity switch. This precision comes at the cost of
recall: the splitter detects only 7.35\% of the 1\,048 ground-truth switches,
producing 99 predicted splits of which 77 are correct. The precision is
expected to be high because jersey-number predictions are only acted upon when
they pass an entropy and persistence filter, ensuring splits are triggered only
when the number change is both confident and temporally sustained. The low
recall reflects the fundamental visibility constraint: jersey numbers are
frequently unreadable due to occlusion, motion blur, and low-resolution crops,
which account for the majority of identity switches in the dataset.

\paragraph{Team splitter.}
The team-colour splitter achieves lower precision (63.51\%) but higher recall
(8.97\%) than the jersey splitter, producing 148 predicted splits of which 94
are correct. Team colour is a coarser signal than jersey number and is
susceptible to lighting variation and partial occlusion, which explains the
increased false-positive rate. However, because team colour can often be
estimated even when jersey numbers are not legible, the team splitter recovers
a slightly broader set of switches.

\paragraph{Jersey and team signals are largely correlated.}
Combining jersey and team yields 154 predicted splits and 96 true positives,
compared to team-only's 148 predicted splits and 94 true positives. The
combination adds only 2 additional true positives while contributing 4
additional false positives, indicating that the two signals detect almost
the same set of switches. This is consistent with the underlying cause: in
most identity-switch events, both the jersey number and the team colour change
simultaneously and are equally (un)observable. The marginal gain of combining
them does not justify added complexity in isolation, but the combined signal
forms a useful base for adding genuinely independent signals.

\paragraph{Bounding-box anomaly splitter.}
The bounding-box anomaly splitter achieves a precision of 61.33\% and a recall
of 8.78\%, closely comparable to the team splitter in both metrics. Despite
similar performance, the bbox signal offers a distinct cost advantage: it
operates purely on the bounding-box trajectory already produced by the tracker
and requires no additional attribute prediction, in contrast to the jersey and
team splitters which require intensive inference pipelines involving pose
estimation, OCR, and embedding extraction. This makes the bbox splitter an
attractive addition to the pipeline — comparable performance to the team
splitter at negligible additional cost. One limitation is that this evaluation
uses ground-truth detections; noisy model detections in practice may introduce
additional spurious velocity spikes that reduce precision. Adding the bbox
splitter to the jersey+team combination increases the true-positive count from
96 to 109, with false positives rising from 58 to 98, and IDS-F1 improving
from 15.97\% to 17.37\%, confirming that the bbox signal provides genuinely
complementary coverage.

\paragraph{Best combination: Jersey + Team + Bbox.}
The jersey, team, and bounding-box signals are unified in a single-pass
splitter that evaluates all three signals on the original, unsplit tracklet.
This avoids the window-truncation artefact of sequential splitting, where the
persistence window of the team signal is artificially clipped at boundaries
introduced by the jersey splitter. The combined splitter achieves the highest
IDS-F1 of all tested configurations (17.37\%), with 109 true positives and 98
false positives across 1\,930 output tracklets.

% ---------------------------------------------------------------------
\subsection{Rejected Alternatives}\label{subsec:rejected_splitters}
% ---------------------------------------------------------------------

Two splitter designs were explored and ultimately rejected from the ALERT
pipeline: the trajectory splitter and the GTA DBSCAN splitter.

\paragraph{Trajectory splitter.}
The trajectory splitter produces 235 predicted splits of which only 82 are
correct, achieving a precision of 34.89\% — the lowest of all methods tested.
The high false-positive rate arises because legitimate football actions —
tackles, close dribbling encounters, and player crossings — produce exactly the
proximity and direction-change pattern the splitter is designed to detect, but
without an actual identity switch. Adding the trajectory splitter on top of the
best combination (jersey+team+bbox) increases the fragment count from 1\,930
to 2\,072 while improving the true-positive count by only 7, at the cost of
135 additional false positives, decreasing IDS-F1 from 17.37\% to 16.61\%.

\paragraph{GTA DBSCAN splitter.}
The GTA DBSCAN splitter is adapted from Sun et al.~\cite{sun2024gta}, who
propose to detect mix-up errors by extracting per-frame box-grained embeddings
using an OSNet ReID model and applying DBSCAN clustering in cosine distance
space. Tracklets whose embeddings form multiple distinct clusters are assumed
to contain multiple identities and are split at the cluster boundaries. The
splitter achieves the highest IDS-Recall of all methods (14.79\%), recovering
switches that the attribute-based signals miss. However, it generates 2\,053
predicted splits for only 155 true positives, yielding a precision of only
7.55\% and a false-positive rate above 92\%. The OSNet model was not trained
on SoccerNet footage, and the resulting domain mismatch causes embeddings to
diverge substantially even within clean tracklets, producing a large number of
spurious cluster boundaries. The resulting fragmentation of 2\,367 output
tracklets from 1\,725 input tracklets — an increase of 37\% — places a very
high burden on the downstream merger.

Nonetheless, the GTA splitter does improve HOTA when followed by the downstream
merger (\texttt{TODO:HOTA:gta}), because its high recall provides the merger
with fragments it would otherwise never see. This suggests that the merger is
sufficiently robust to absorb the false-positive splits while benefiting from
the additional true positives. The result motivates future work on
domain-adapted ReID models that could combine GTA's recall with higher
precision.

% ---------------------------------------------------------------------
\subsection{Unified vs.\ Sequential Splitting}\label{subsec:unified_vs_sequential}
% ---------------------------------------------------------------------

The jersey, team, and bounding-box signals can be applied either sequentially
— where each splitter receives the already-split output of the previous one —
or in a unified single pass over the original tracklet. In the sequential
design, the team splitter operates on tracklets that have already been split
by the jersey splitter, so its persistence lookahead window is clipped at the
jersey-introduced fragment boundaries. In the unified design, all three
signals evaluate the original tracklet simultaneously and their switch indices
are merged before a single split is applied, preserving the full lookahead
context for every signal.

The difference is measurable: the unified jersey+team splitter achieves an
end-to-end HOTA of \texttt{TODO:HOTA:unified} compared to
\texttt{TODO:HOTA:sequential} for the sequential design, a gain of
approximately 0.007 HOTA points. The improvement is attributable to the
unified design suppressing false splits near jersey-split boundaries, where
the sequential team splitter sees a truncated persistence window and
incorrectly fires on short spurious team-signal segments that span the
boundary.

% ---------------------------------------------------------------------
\subsection{Summary}\label{subsec:splitter_summary}
% ---------------------------------------------------------------------

The jersey-number signal is the most precise individual splitter (77.78\%)
but covers fewer than 8\% of ground-truth switches, reflecting the limited
visibility of jersey numbers in broadcast footage. The team-colour and
bounding-box anomaly signals achieve comparable precision (61--64\%) and
slightly higher recall, and are largely orthogonal to the jersey signal; the
bbox signal additionally contributes at negligible computational cost relative
to the attribute-based signals. Combining all three in a unified single-pass
splitter yields the best intrinsic performance (IDS-F1 = 17.37\%) and is
adopted as the ALERT splitter. The trajectory and GTA DBSCAN splitters both
fail on precision — the former because football motion inherently mimics the
proximity-swap pattern, the latter because the OSNet ReID model was not trained on SoccerNet footage,
causing embeddings to diverge even within clean tracklets — and are excluded
from the final pipeline. Across all
configurations, IDS-Recall remains below 15\%, reflecting that the majority of
identity switches occur during occlusion or at distances where no signal is
reliably observable.
